\chapter{Introducción}
\label{chap:introduccion}

% Versión 2 del capítulo 1 (2026-10-04). Estructura indicada por la dirección: contexto y
% problema, justificación, objetivos (texto literal del plan de trabajo de grado), alcance y
% organización. No adelanta resultados; es coherente con los capítulos 2 a 5.

Conocer la distribución de densidad en el interior de un edificio volcánico es una de las bases para evaluar su amenaza. Los conductos por los que asciende el magma, las cavidades, las zonas de alteración hidrotermal y los cuerpos de roca más densa modifican la masa que encuentra cualquier trayecto que atraviese la montaña, y su localización contribuye a interpretar procesos como la desgasificación, el ascenso de material o la estabilidad de los flancos \parencite{tanaka2007asama,tanaka2023muography}. Los métodos geofísicos convencionales infieren esa estructura de manera indirecta, a partir de campos medidos en la superficie, y las técnicas de inspección con radiación artificial no alcanzan a atravesar cientos de metros de roca \parencite{bonomi2020}. Este trabajo se ocupó de una técnica que sí lo hace, la muografía por transmisión, y de una de sus limitaciones más severas: el tiempo de exposición necesario para obtener una imagen interpretable. El capítulo presenta el problema que motivó el estudio, las razones que justificaron abordarlo con aprendizaje profundo, los objetivos formulados, el alcance del trabajo y la organización del documento.

\section{Contexto y planteamiento del problema}
\label{sec:intro-problema}

La muografía aprovecha los muones que se producen cuando los rayos cósmicos interactúan con la alta atmósfera. Por su elevada energía y su capacidad de penetración, una fracción de esos muones atraviesa cientos de metros, e incluso kilómetros, de roca, y la atenuación de su flujo en cada dirección de observación informa sobre la opacidad, es decir, sobre la cantidad de materia que encuentran a lo largo de ella \parencite{bonomi2020,tanaka2023muography}. En la modalidad por transmisión, un detector situado al pie del edificio registra cuántos muones llegan desde cada dirección, y el resultado es un mapa bidimensional en el que cada píxel resume la densidad acumulada a lo largo de un trayecto. Con esta técnica se visualizaron estructuras bajo el cráter del monte Asama \parencite{tanaka2007asama} y se obtuvieron imágenes de alta definición del volcán Sakurajima \parencite{olah2018sakurajima}, ambos en Japón, y en Colombia se diseñó y construyó el telescopio híbrido MuTe con el propósito de estudiar volcanes activos del país \parencite{penarodriguez2020}.

La fuente de la técnica es, al mismo tiempo, su principal limitación. A diferencia de la radiografía médica o industrial, en la que la intensidad de la fuente puede ajustarse, la muografía depende del flujo natural de muones, cercano a un muon por centímetro cuadrado y por minuto sobre un detector horizontal al nivel del mar \parencite{pdg2022}. En las direcciones casi horizontales desde las que se observa un volcán ese flujo es mucho menor, y después de atravesar cientos de metros de roca solo sobrevive una fracción pequeña de él. La llegada de los muones es, además, un proceso aleatorio de conteo: el número registrado en cada dirección fluctúa alrededor de su valor esperado, y el ruido relativo de esa fluctuación disminuye solo con la raíz cuadrada del número de eventos, como se expone en la \cref{sec:mt-estadistica}. Para que los contrastes de densidad se distingan de las fluctuaciones se requieren exposiciones de semanas o meses \parencite{lesparre2010}, durante las cuales el instrumento debe operar en ambientes remotos, a la intemperie y con suministro de energía limitado. Esa condición encarece la adquisición, aumenta el riesgo de fallas y limita la frecuencia con la que puede actualizarse una imagen \parencite{bonomi2020,lechmann2021}.

Una exposición corta produce, en consecuencia, una imagen de baja estadística. En ella, la fluctuación de conteo domina sobre las diferencias que se pretende observar, las direcciones de menor transmisión registran pocas cuentas o ninguna, y el contraste de una anomalía de densidad puede quedar enmascarado o confundirse con una fluctuación del mismo orden. Un suavizado de la imagen reduce el ruido, pero puede atenuar también los contrastes débiles y las estructuras de pequeña escala que sustentan la interpretación. Disminuir el tiempo de adquisición sin perder la información necesaria para interpretar la imagen constituye, por ello, un problema abierto de la muografía volcánica.

En otros campos de la imagenología, el aprendizaje profundo ha ofrecido una alternativa al aumento de la exposición. En la tomografía computarizada de baja dosis, cuyo ruido cuántico también sigue una distribución de Poisson, las redes generativas antagónicas se han empleado para transformar imágenes ruidosas en estimaciones próximas a las de dosis completa \parencite{yang2018}. En muografía, \textcite{odonnell2025particles} aplicaron una red generativa antagónica condicional de Wasserstein con penalización de gradiente a imágenes de baja estadística de bloques de concreto y reportaron que imágenes de un día de exposición alcanzaron, según el índice de similitud estructural, la calidad de imágenes de 21 días. Ese resultado se obtuvo, sin embargo, en un escenario de infraestructura civil, cuyas geometrías, distancias y contrastes de densidad difieren sustancialmente de los de un edificio volcánico.

De esta situación surgió la pregunta que orientó el trabajo: si un modelo de aprendizaje profundo, entrenado con datos sintéticos generados mediante simulación Monte Carlo, puede transformar imágenes muográficas por transmisión de baja estadística en imágenes de calidad comparable a las obtenidas con exposiciones prolongadas en escenarios volcánicos, conservando los contrastes de densidad que dan sentido a la imagen, y en qué medida ese procedimiento permitiría reducir el tiempo de adquisición. La reducción del tiempo de adquisición fue, así, el propósito investigado; lo que la evidencia obtenida permite afirmar al respecto se expone en el \cref{chap:conclusiones}.

\section{Justificación}
\label{sec:intro-justificacion}

El estudio se justificó por tres razones: la pertinencia del aprendizaje profundo para el problema planteado, la necesidad de evaluarlo con criterios propios de la muografía y la viabilidad de abordarlo mediante simulación.

La primera razón es el costo del tiempo de exposición y la posibilidad de reducirlo por medios computacionales. Como disminuir a la mitad el ruido relativo de un conteo exige cuadruplicar la exposición, cualquier procedimiento que extraiga más información de una adquisición corta incide directamente sobre el costo logístico y sobre la frecuencia con la que puede vigilarse un volcán \parencite{tanaka2023muography,lechmann2021}. Un modelo supervisado puede aprender, a partir de pares de imágenes de baja y alta estadística, regularidades espaciales asociadas al proceso de atenuación y estimar con ellas la imagen que se habría obtenido con mayor exposición. Las redes neuronales se han aplicado ya a la reducción del ruido estadístico de muografías simuladas \parencite{lefevre2024} y a la tomografía por dispersión de corta duración \parencite{wang2026}, y el trabajo de \textcite{odonnell2025particles} mostró la factibilidad de la mejora predictiva en estructuras de concreto. En la literatura revisada, que se presenta en la \cref{sec:mt-estado-arte}, no se identificó, en cambio, la aplicación de redes generativas condicionales a imágenes de baja estadística de muografía volcánica por transmisión.

La segunda razón es que la utilidad de una imagen mejorada no puede juzgarse solo por su parecido con la referencia. En la muografía, cada píxel es un conteo ligado a la opacidad del trayecto correspondiente, y la interpretación se apoya en las cuentas de regiones y en los contrastes entre ellas. Las métricas de similitud de uso general, como el índice de similitud estructural y la relación señal-ruido de pico, resumen la imagen completa y no distinguen las regiones relevantes para la interpretación. Además, la distorsión respecto de la referencia y la apariencia realista de una estimación no pueden optimizarse simultáneamente más allá de cierto límite \parencite{blau2018}, y en la traducción de imágenes médicas se ha mostrado que las pérdidas que ajustan la distribución de las imágenes generadas pueden introducir o suprimir estructuras ausentes en la imagen real \parencite{cohen2018}. Una red que suaviza el ruido puede, por tanto, elevar la similitud global y, al mismo tiempo, desplazar cuentas o atenuar una anomalía. Por ello se consideró necesario evaluar el esquema tanto por la similitud de sus imágenes con la referencia como por la conservación de la información relevante para la muografía, es decir, las cuentas regionales, el contraste de las anomalías y su signo, con los criterios que se fundamentan en la \cref{sec:mt-metricas}.

La tercera razón es la viabilidad del estudio mediante simulación. En campo pueden obtenerse pares de imágenes de corta y larga exposición; sin embargo, la referencia conserva fluctuaciones estadísticas y la estructura interna verdadera del objeto generalmente no se conoce. La simulación permite construir estos pares sobre escenarios de geometría y materiales conocidos. La simulación Monte Carlo del transporte de muones con Geant4 \parencite{agostinelli2003,allison2016}, alimentada por la parametrización del flujo de \textcite{guan2015muonflux} y por el generador EcoMug \parencite{pagano2021}, permite construir esos pares en condiciones controladas, sobre escenarios cuya estructura se fija de antemano y sin depender de campañas experimentales costosas. Esa posibilidad hizo abordable el problema en el marco de un trabajo de grado y facilita la reproducción de sus resultados. El proyecto se inscribió, además, en la línea de trabajo del Grupo de Óptica y Tratamiento de Señales de la Universidad Industrial de Santander, que articula la simulación física, el procesamiento de imágenes y el aprendizaje profundo, y en el interés nacional por la muografía de volcanes colombianos \parencite{penarodriguez2020}.

\section{Objetivos}
\label{sec:intro-objetivos}

Los objetivos se presentan a continuación tal como se formularon en el plan de trabajo de grado.

\subsection{Objetivo general}
\label{subsec:intro-objetivo-general}

Desarrollar y validar un esquema de reconstrucción muográfica por transmisión asistido por aprendizaje profundo para reducir el tiempo efectivo de adquisición en exploración volcánica, mediante mejora predictiva de imágenes de baja estadística hacia imágenes de referencia de alta estadística, evaluando la mejora en calidad y detectabilidad de estructuras internas.

\subsection{Objetivos específicos}
\label{subsec:intro-objetivos-especificos}

\begin{enumerate}[label=\textbf{OE\arabic*.}, leftmargin=*, labelsep=0.6em]
  \item Generar un conjunto de datos sintéticos de imágenes muográficas por transmisión, con diferentes niveles de estadística de conteo, mediante simulación Monte Carlo que reproduzca escenarios representativos de exploración volcánica.
  \item Diseñar e implementar un modelo de aprendizaje profundo basado en redes generativas adversarias condicionales para transformar imágenes muográficas de baja estadística en imágenes de calidad comparable a las de alta estadística.
  \item Evaluar el desempeño del modelo entrenado y determinar el factor de reducción del tiempo de adquisición alcanzable con el esquema propuesto, comparando la calidad de las imágenes mejoradas por el modelo frente a imágenes obtenidas con tiempos de exposición prolongados.
\end{enumerate}

Los tres objetivos específicos se abordaron en secuencia, porque el producto de cada uno fue el insumo del siguiente: el conjunto sintético del primero sirvió para entrenar y evaluar el modelo del segundo, y ese modelo fue el objeto de la evaluación del tercero. El objetivo general añade a esa secuencia la exigencia de juzgar las imágenes estimadas tanto por su calidad como por la detectabilidad de las estructuras internas.

\section{Alcance del trabajo}
\label{sec:intro-alcance}

El plan de trabajo de grado delimitó el estudio a geometrías simplificadas, manejables computacionalmente a nivel de pregrado, y a mapas bidimensionales obtenidos a partir de la atenuación del flujo de muones. Dentro de ese marco, el trabajo fue un estudio enteramente computacional con tres componentes: la definición de escenarios volcánicos sintéticos, la simulación Monte Carlo del transporte de muones a través de ellos y la evaluación de modelos de aprendizaje profundo sobre las imágenes resultantes.

Los escenarios fueron conos idealizados de roca homogénea con anomalías internas de geometría simple, del tipo que la muografía volcánica asocia con conductos, cavidades y cuerpos densos. No reproducen la topografía ni la estructura de densidades de un volcán concreto, y su representatividad es de mecanismos y configuraciones, no de una población geológica. Cada imagen es un mapa de conteos direccionales, ordenado por elevación y acimut desde una posición de observación, que registra los muones primarios que alcanzan un plano de conteo con la energía establecida. Las imágenes no incluyen la respuesta de un detector físico, el fondo instrumental ni las partículas secundarias, y su exposición es una escala convencional y no un tiempo calibrado de medición. En consecuencia, el trabajo no incluyó adquisición de datos de campo, reconstrucción tridimensional ni una inversión de opacidad validada: la información de densidad de las imágenes estimadas se examinó mediante contrastes de cuentas, que no equivalen a diferencias de opacidad (\cref{sec:mt-imagen}).

Por la misma razón, la reducción del tiempo de adquisición se estudió bajo una equivalencia idealizada. Con flujo estable y detección constante, el número esperado de cuentas es proporcional a la exposición, de modo que reducir los conteos de una imagen de referencia mediante submuestreo equivale nominalmente a reducir su tiempo de exposición. Toda razón de reducción que se discute en el documento es, por ello, una razón de conteos en condiciones simuladas, y su traducción a un tiempo de adquisición de campo requeriría un instrumento calibrado.

En el componente de aprendizaje, los modelos se entrenaron con una sola fracción de conteo y se evaluaron después con varias, con el fin de examinar su comportamiento ante distintas estadísticas de entrada. Los escenarios se asignaron por geometría completa a entrenamiento, validación, prueba y generalización, con doce, tres, tres y una geometrías, respectivamente, en lugar de la proporción de 70, 15 y 15\,\% prevista en el plan de trabajo. La asignación por geometría impide que imágenes derivadas de un mismo escenario intervengan a la vez en el ajuste y en la evaluación \parencite{kaufman2012}, a costa de un número reducido de escenarios reservados. Los indicadores de conservación de anomalías se calcularon sobre regiones cuya ubicación se conocía por construcción; permiten examinar si una estructura conocida se preserva, pero no constituyen un procedimiento de detección ni de localización de anomalías desconocidas.

Las afirmaciones del trabajo se refieren a estas condiciones. La inversión de las imágenes estimadas a opacidad y la validación del esquema con un instrumento real en campo quedaron fuera del estudio; su desarrollo permitiría establecer si lo observado en el caso sintético se sostiene ante observaciones reales, y se plantea como posible continuación en la \cref{sec:trabajo-futuro}.

\section{Organización del documento}
\label{sec:intro-organizacion}

El documento se organiza en cinco capítulos. Después de esta introducción, el \cref{chap:marco-teorico} presenta los fundamentos en que se apoya el trabajo: el paso de los muones por la materia, la formación de la imagen muográfica por transmisión, la estadística de conteo, la simulación Monte Carlo, los modelos de aprendizaje profundo para la mejora de imágenes, las métricas de calidad y conservación y los antecedentes del aprendizaje profundo en muografía. El \cref{chap:metodologia} describe el diseño del estudio, la simulación de los escenarios, la verificación y preparación del conjunto de datos, la arquitectura y el entrenamiento de los modelos y los protocolos de evaluación. El \cref{chap:resultados} presenta y discute los resultados de cada objetivo específico, y el \cref{chap:conclusiones} expone las conclusiones por objetivo, la conclusión general y posibles líneas de continuación. El anexo~\ref{chap:anexo-protocolo} detalla un protocolo propuesto, no ejecutado, para estimar el factor de reducción con criterios fijados de antemano.
